\documentclass[11pt,a4paper]{article}

% USE LINUX LIBERTINE FONTS
\usepackage{libertine}
\usepackage{libertinust1math}
% AND INCONSOLATA FOR MONOSPACED
\usepackage{inconsolata}
%% PACHETE MATEMATICE
%\usepackage{amsmath,amssymb,amsfonts}
\usepackage{amsthm}
% \usepackage{mathpazo}
% \usepackage{eulervm}
\usepackage{enumerate}
\usepackage{color}
\usepackage{graphicx}
\usepackage[all]{xy}
\usepackage{xcolor}
\usepackage{framed}
\usepackage[unicode]{hyperref}
\setcounter{MaxMatrixCols}{20}
\usepackage{multicol}
%\usepackage[romanian]{babel}


%% ASPECT
\linespread{1.05}
\usepackage[T1]{fontenc}
\usepackage[utf8x]{inputenc}
\usepackage[margin=2cm]{geometry}
% \usepackage[color]{showkeys}
\usepackage{anyfontsize}
\usepackage{titlesec}
\titleformat*{\section}{\large\bfseries}

\usepackage{fancyhdr}
\pagestyle{fancy}
\rhead{\small \color{gray} Notițe de Adrian Manea}
\lhead{\small \color{gray} ETTI-AG, 2017---2018, A. Niță \& A. Oprina}


\usepackage[autostyle = false, style = german]{csquotes}
\usepackage[romanian]{babel}
\newcommand\qq{\enquote}

%% COMENZILE MELE
\renewcommand*{\proofname}{Dem.:}
\newcommand\bb{\mathbb}
\newcommand\dr{\mathrm}
\newcommand\kal{\mathcal}
\newcommand\fk{\mathfrak}
\newcommand\op{\oplus}
\newcommand\ot{\otimes}
\newcommand\ds{\displaystyle}
\newcommand\id{\indent}
\newcommand\nid{\noindent}
\newcommand\seq{\subseteq}
\newcommand\sneq{\subsetneq}
\newcommand\speq{\supseteq}
\newcommand\spneq{\supsetneq}
\newcommand\rar{\rightarrow}
\newcommand\Rar{\Rightarrow}
\newcommand\xrar{\xrightarrow}
\newcommand\lar{\leftarrow}
\newcommand\Lar{\Leftarrow}
\newcommand\xlar{\xleftarrow}
\newcommand\lrar{\leftrightarrow}
\newcommand\Lrar{\Leftrightarrow}
\newcommand\ideal{\trianglelefteq}
\newcommand\ai{\text{ a.\^{i}. }}
\newcommand\sk{(\textit{Schi\c{t}\u{a}}) }
\newcommand\rhup{\rightharpoonup}
\newcommand\rhdn{\rightharpoondown}
\newcommand\lhup{\leftharpoonup}
\newcommand\lhdn{\leftharpoondown}
\newcommand\vid{\emptyset}
\newcommand\wh{\widehat}
\newcommand\ol{\overline}
\newcommand\NN{\mathbb{N}}
\newcommand\RR{\mathbb{R}}
\newcommand\ZZ{\mathbb{Z}}
\newcommand\QQ{\mathbb{Q}}
\newcommand\CC{\mathbb{C}}
\newcommand\Bbbk{\mathbb{k}}

%% STILURILE MELE
\newtheoremstyle{dotless}{}{}{\itshape}{}{\bfseries}{:}{ }{}

\theoremstyle{dotless}
\newtheorem{theorem}{Teorem\u{a}}[section]
\newtheorem{proposition}{Propozi\c{t}ie}[section]
\newtheorem{lemma}{Lem\u{a}}[section]
\newtheorem{corollary}{Corolar}[section]
\newtheorem{exercise}{Exerci\c{t}iu}

\newtheoremstyle{dotlessdr}{}{}{}{}{\bfseries}{:}{ }{}
\theoremstyle{dotlessdr}
\newtheorem{example}{Exemplu}[section]
\newtheorem{definition}{Defini\c{t}ie}[section]
\newtheorem{remark}{Observa\c{t}ie}[section]




\begin{document}

\begin{center}
  {\large\textbf{Probleme propuse de Prof. Oprina }}
\end{center}

\vspace{1cm}


1. Definim mulțimea:
\[
  \CC_n[X] = \{ P \in \CC[X] \mid \mathrm{grad} P \leq n \}.
\]

Arătați că:
\begin{enumerate}[(a)]
\item $(\CC_n[X], +)$ este un $\CC$-spațiu vectorial de dimensiune $n + 1$;
\item $\forall a \in \CC$ fixat, mulțimea:
  \[
    B = \{1, X - a, (X - a)^2, \dots, (X - a)^n \}
  \]
  este o bază în $\CC_n[X]$;
\item Dacă $f \in \CC_n[X]$, atunci:
  \[
    f(x) = \sum_{k = 0}^n \frac{f^{(k)}(x)}{k!} (X - a)^k \quad (\text{Taylor}).
  \]
\end{enumerate}

\vspace{1cm}

2. Fie $\Bbbk$ un corp comutativ și considerăm sistemul liniar $AX = 0$,
cu $A \in M_n(\Bbbk)$.

Arătați că:
\begin{enumerate}[(a)]
\item $S \leq \Bbbk^n$, unde $S$ este mulțimea soluțiilor sistemului;
\item Există $\varphi : \Bbbk^n \to \Bbbk^n$ liniară astfel încît $\mathrm{Ker} \varphi = S$ (Indicație: $\varphi(x) = Ax$);
\item $\dim_\Bbbk S = n - r$, unde $R = \mathrm{rang}(A)$.
\end{enumerate}

\vspace{1cm}

3. Fie $A \in M_n(\RR)$, cu $A \neq 0$ și $\det A = 0$.

Să se arate că există $B \in M_n(\RR)$, cu $B \neq 0$, astfel încît $AB = 0$.

\end{document}
